(1) 内部无旋且不可压，\(\nabla^2\phi=0\)。正则 \(\ell\) 阶球谐解只能选 \(r^\ell Y_{\ell m}\)，中心奇异的 \(r^{-\ell-1}\) 项被排除。线性运动学边界 \(v_r(R)=\dot\eta Y_{\ell m}\)，故
\[
\boxed{\phi(r,\theta,\varphi,t)=\frac{\dot\eta(t)}{\ell R^{\ell-1}}r^\ell Y_{\ell m}}.
\]
球内速度为 \(\nabla\phi\)，表面满足 \(v_r=\dot\eta Y\)。该模型以真空为外界，未计任何外流体的附加惯量。

(2) 球面平均曲率的一阶扰动为 \(\delta(2H)=(\ell-1)(\ell+2)\eta Y/R^2\)，由Laplace压差得 \(\delta p=\sigma(\ell-1)(\ell+2)\eta Y/R^2\)。线性非定常Bernoulli在表面给 \(\delta p=-\rho\,\partial_t\phi|_R=-\rho R\ddot\eta Y/\ell\)。两式相等得到
\[
\boxed{\ddot\eta+\omega_\ell^2\eta=0,\quad
\omega_\ell^2=\frac{\sigma\ell(\ell-1)(\ell+2)}{\rho R^3}}.
\]
频率仅依赖 \(\ell\) 不依赖 \(m\)，故在球对称自由液滴中每个 \(\ell\) 有 \(2\ell+1\) 重简并；外部约束会打破此简并。

(3) 体动能可由Green恒等式求：
\(T=\frac\rho2\int_{r<R}|\nabla\phi|^2dV=\frac\rho2\oint_R\phi\partial_r\phi\,dS
=\frac{\rho R^3}{2\ell}|\dot\eta|^2\)，使用了单位归一球谐。表面张力在固定总体积条件下的二阶面积变化
\(\Delta A=\frac12(\ell-1)(\ell+2)|\eta|^2\)，故
\[
V_\sigma=\frac{\sigma}{2}(\ell-1)(\ell+2)|\eta|^2.
\]
广义质量 \(M_\ell=\rho R^3/\ell\)、刚度 \(K_\ell=\sigma(\ell-1)(\ell+2)\)，比值 \(K_\ell/M_\ell\) 正是上问 \(\omega_\ell^2\)，构成独立的能量法核验。

(4) \(\ell=1\) 对应将球整体平移而不改变形状，频率为零；\(\ell=0\) 对应总体积变化，若液体严格不可压且总质量固定，它不属于可自由振荡的形变自由度。固定 \(\ell,\sigma,\rho\) 时 \(\omega_\ell\propto R^{-3/2}\)，微小液滴频率增大；但当 \(R\) 接近分子尺度时连续介质与恒定表面张力近似失效。令 \(\sigma\to0\) 时无形变恢复力，所有该类毛细模趋于零。
